% vim: ai sw=2 spell spelllang=cs

\begin{frame}{Programování v síti}

  \begin{itemize}
    \item Komunikace mezi vzdálenými stroji
    \begin{itemize}
      \item Ale i lokálně
    \end{itemize}

    \pause

    \item Stroje musí mít vzor, podle kterého komunikovat
    \begin{itemize}
      \item Tzv.\ protokol
    \end{itemize}

    \pause

    \item Spousta protokolů
    \begin{itemize}
      \item ARP, RARP
      \item \emph{IPv4, IPv6}
      \begin{itemize}
	\item DHCP, ICMP, \emph{TCP}, UDP
      \end{itemize}

      \item IPsec
      \begin{itemize}
	\item IPv4, IPv6
      \end{itemize}

      \item IPX
      \begin{itemize}
	\item SPX
      \end{itemize}
    \end{itemize}
  \end{itemize}

\end{frame}

\ifpcdirlinuxb

\begin{frame}{Vlastnosti}

  \begin{itemize}
    \item X
  \end{itemize}

\end{frame}

% vim: ai sw=2 spell spelllang=cs

\begin{frame}{UNIX a síť}

  \begin{itemize}
    \item Na UNIXu býval všechno soubor

    \pause

    \item Na soubor existuje souborový deskriptor

    \pause

    \item I na síťové spojení se vytváří souborový deskriptor
    \begin{itemize}
      \item Tzv.\ socket (\doc{man 7 socket})
    \end{itemize}

    \pause

    \item Na něm se dá provádět většina standardních operací
  \end{itemize}

\end{frame}

\begin{frame}[fragile]{\code{socket}}

  \begin{itemize}
    \item Otevření: \code{int socket(int domain, int type, int protocol)}
    \begin{itemize}
      \item Hlavičky: \header{sys/types.h}, \header{sys/socket.h}
      \item \code{domain}: protokol -- \code{AF_UNIX} \code{AF_INET},
	\code{AF_INET6}, \dots
      \item \code{type}: \code{SOCK_STREAM} (TCP), \code{SOCK_DGRAM} (UDP), \dots
      \item \code{protocol}: pokud je na výběr (NETLINK), jinak 0 (pro UDP i TCP)
      \item Návratová hodnota: souborový deskriptor nebo \code{\-1} při chybě
    \end{itemize}

    \pause

    \item Zavření
    \begin{itemize}
      \item Klasicky: \code{close}
\ifpcdirlinuxb
      \item Částečně: \code{int shutdown(int sockfd, int how)}
      \begin{itemize}
	\item \code{how} -- \code{SHUT_RD}, \code{SHUT_WR}, \code{SHUT_RDWR}
      \end{itemize}
\fi
    \end{itemize}

    \pause
  \end{itemize}

  \onslide<1->

  \begin{block}{Příklad (bez ošetření chyb)}
  \begin{lstlisting}
int s = socket(AF_UNIX, SOCK_STREAM, 0);
close(s);
  \end{lstlisting}
  \end{block}

\end{frame}


\begin{frame}{Operace nad sockety}

  \begin{itemize}
    \item Komunikace
    \begin{itemize}
      \item Standardně (\code{read}, \code{write}, \code{readv},
	\code{writev}, \code{sendfile})
      \item Síťovými prostředky (\code{recv}, \code{send}, \dots)
    \end{itemize}

    \pause

    \item Multiplexing
    \begin{itemize}
      \item Standardně (\code{select}, \code{poll})
    \end{itemize}

    \pause

    \item Parametry
    \begin{itemize}
      \item \code{ioctl}, \code{fcntl} (\code{O_NONBLOCK})
      \item \code{getsockopt}, \code{setsockopt}
    \end{itemize}

    \pause

    \item Spojování
    \begin{itemize}
      \item Klient: \code{connect}
      \item Server: \code{bind}, \code{listen}, \code{accept}
    \end{itemize}
  \end{itemize}

\end{frame}

\begin{frame}[fragile]{Síťové prostředky komunikace}

  \begin{block}{Prototypy}
  \begin{lstlisting}[aboveskip=2pt,belowskip=3pt]
ssize_t recv(int sock, void *buf, size_t len, int flags)
ssize_t send(int sock, const void *buf, size_t len, int flags)
  \end{lstlisting}

  \begin{onslide}<2->
    \begin{lstlisting}[aboveskip=3pt,belowskip=2pt]
ssize_t recvfrom(int sock, void *buf, size_t len, int flags,
      struct sockaddr *src_addr, socklen_t *addrlen)
ssize_t sendto(int sock, const void *buf, size_t len, int flags,
      const struct sockaddr *dest_addr, socklen_t addrlen)
  \end{lstlisting}
  \end{onslide}
  \end{block}

  \begin{itemize}
    \item Lze kontrolovat odeslání a příjem
    \begin{itemize}
      \item \code{flags} -- \code{MSG_DONTWAIT}, \code{MSG_OOB},
	\code{MSG_PEEK}, \code{MSG_DONTROUTE}
    \end{itemize}

    \pause

    \item Lze adresovat
    \begin{itemize}
      \item Jen pro nespojované sockety (dále)
    \end{itemize}
  \end{itemize}

  \pause

  \begin{block}{Příklad (bez ošetření chyb)}
  \begin{lstlisting}[aboveskip=2pt,belowskip=2pt]
int sock = socket(AF_INET, SOCK_DGRAM, 0);
sendto(sock, "ahoj", 4, addr, sizeof(addr));
  \end{lstlisting}
  \end{block}

\end{frame}


\subsection*{Unix sockety}
\frame{\subsectionpage}

\begin{frame}[fragile]{Unix sockety}

  \begin{itemize}
    \item Adresace cestou v souborovém systému
    \begin{itemize}
      \item Nebo bezejmenný
    \end{itemize}

    \pause

    \item Jen lokální komunikace

    \pause

    \item Protokol: \code{AF_UNIX} (\doc{man 7 unix})

    \pause

    \item Vytvoření
    \begin{itemize}
      \item \code{socketpair} (už známe)
      \item \code{socket} (viz dále)
    \end{itemize}
  \end{itemize}

  \onslide<4->

  \begin{block}{Příklad (bez ošetření chyb)}
  \begin{lstlisting}
socketpair(AF_UNIX, SOCK_STREAM, 0, sock);
write(sock[0], "ahoj", 4);
read(sock[1], buf, 4);
  \end{lstlisting}
  \end{block}

\end{frame}

\begin{frame}{Úkol}

  \heading{Komunikace přes unixový socket}
  \begin{enumerate}
    \item Otevřete si pár socketů (\code{socketpair})
    \item Zapište do 1 socketu postupně 2 řetězce
    \begin{itemize}
      \item 2 $\times$ \code{write} se stejným socketem
    \end{itemize}

    \item Přečtěte z 2.\ socketu data
    \item Přeložte, vyzkoušejte

    \item Kolikrát je třeba číst na přečtení 2 zápisů?
    \begin{itemize}
      \item Vysvětlete
    \end{itemize}
  \end{enumerate}

\end{frame}

\subsection*{Protokol IP}
\frame{\subsectionpage}

\fi

\begin{frame}{Protokol IP}

  \begin{itemize}
    \item Internet -- síť s přepojováním paketů pomocí IP

    \pause

    \item Směrovače -- spojují více sítí
    \begin{itemize}
      \item Směrování podle cílové sítě
    \end{itemize}

    \pause

    \item Síť je homogenní, používá stejné protokoly

    \pause

    \item Dnes dvě verze
    \begin{enumerate}
      \item IPv4
      \begin{itemize}
	\item 32bitové adresy
	\item \doc{man 7 ip}
      \end{itemize}

      \item IPv6
      \begin{itemize}
	\item 128bitové adresy
	\item \doc{man 7 ipv6}
      \end{itemize}
    \end{enumerate}
  \end{itemize}

\end{frame}

\begin{frame}{Adresace v IP}

  \heading{IPv4}
  \begin{itemize}
    \item 32bitová, udává se jako 4 desítková čísla oddělená tečkou
    \item Za lomítkem může být maska sítě

    \pause

    \item Např.\ \cmd{10.0.0.1/24}
    \begin{itemize}
      \item Adresa \cmd{10.0.0.1}
      \item Patří do sítě \cmd{10.0.0.255}
      \begin{itemize}
	\item Všichni v okolí mají adresu \cmd{10.0.0.X}
	\item Nemůže přímo komunikovat s nikým mimo svoji síť
      \end{itemize}
    \end{itemize}
  \end{itemize}

  \heading{IPv6}
  \begin{itemize}
    \item 128bitová, udává se jako dvojice šestnáctkových čísel
      oddělených dvojtečkou
    \begin{itemize}
      \item Nuly lze vynechat
    \end{itemize}

    \pause

    \item Např.\ \cmd{2001:718:801:230::1/64}
  \end{itemize}

  \pause

  \begin{center}
  \heading{Kdo si má adresu pamatovat?} \\
  \pause
  \heading{DNS}
  \end{center}

\end{frame}

\begin{frame}{Protokoly nad IP}

  \heading{Nad IP nejčastěji}
  \begin{itemize}
    \item UDP -- nespojované
    \begin{itemize}
      \item Málo garancí -- nespolehlivá
      \item Bez latencí

      \pause

      \item Výměna jednotlivých datagramů

      \pause

      \item Vhodná na zvuk, video, hry, kde výpadky nevadí
    \end{itemize}

    \pause

    \item TCP -- spojované
    \begin{itemize}
      \item Garantuje pořadí paketů, ověřuje doručení -- spolehlivá
      \item Ale vyšší latence

      \pause

      \item Výměna proudu bajtů

      \pause

      \item Vhodná na data
    \end{itemize}

    \pause

    \item Jedna IP může obsluhovat více UDP i TCP služeb
    \begin{itemize}
      \item Obě podporují tzv.\ porty
      \item 16bitové číslo -- 64k služeb
    \end{itemize}
  \end{itemize}

\end{frame}

\ifpcdirlinuxb

\subsection*{UDP}
\frame{\subsectionpage}

\begin{frame}{Protokol UDP/IP}

  \heading{UDP komunikace}
  \begin{itemize}
    \item Server poslouchá
    \begin{itemize}
      \item Předem daný UDP port (\code{bind} dále)
      \item Nebo náhodný, volný

      \pause

      \item Provede se \code{read}, \code{recvfrom}, \dots
    \end{itemize}

    \pause

    \item Klient pošle datagram
    \begin{itemize}
      \item Musí vědět/umět zjistit IP adresu i UDP port

      \pause

      \item Provede se \code{write}, \code{sendto}, \dots
    \end{itemize}

    \pause

    \item Samostatný datagram ne/proletí sítí
    \begin{itemize}
      \item Další může ne/projít jinudy, rychleji, před prvním apod.
    \end{itemize}
  \end{itemize}

\end{frame}

\begin{frame}{Úkol}

  \heading{UDP komunikace}
  \begin{enumerate}
    \item Otevřete UDP socket na IPv4
    \begin{itemize}
      \item \code{socket}, \code{AF_INET} a \code{SOCK_DGRAM}
    \end{itemize}

    \item Proveďte do něj \code{write} prázdného řetězce a délky 0
    \begin{itemize}
      \item Selže, ale socketu přidělí UDP port
    \end{itemize}

    \item Udělejte \code{recv} do bufferu
    \item Buffer vypište
    \item Zavřete socket
    \item Přeložte a spusťte

    \item Zjistěte port
    \begin{itemize}
      \item Pomocí \cmd{ss -ulpn}
    \end{itemize}

    \item Pošlete něco na UDP port
    \begin{itemize}
      \item Pomocí \cmd{netcat -u 127.0.0.1 port}
    \end{itemize}
  \end{enumerate}

\end{frame}

\subsection*{TCP}
\frame{\subsectionpage}

\fi

\begin{frame}{Protokol TCP/IP}

  \heading{TCP spojení}
  \begin{itemize}
    \item Server poslouchá
    \begin{itemize}
      \item Předem daný TCP port
      \item Nebo náhodný, volný
    \end{itemize}

    \pause

    \item Klient se připojí
    \begin{itemize}
      \item Musí vědět/umět zjistit IP adresu i TCP port
    \end{itemize}

    \pause

    \item Vznikne čtveřice a trvá po celé spojení
    \begin{itemize}
      \item Zdrojová IP a TCP port
      \item Cílová IP a TCP port
    \end{itemize}

    \pause

    \item Proud dat

    \pause

    \begin{itemize}
      \item Pakety IP se mohou zpozdit, nedoručit, předběhnout
      \item TCP korektně doručí všechny a ve správném pořadí
      \begin{itemize}
	\item Může trvat dlouho
      \end{itemize}
    \end{itemize}
  \end{itemize}

\end{frame}

\ifpcdirlinux
% vim: ai sw=2 spell spelllang=cs

\begin{frame}{UNIX a síť}

  \begin{itemize}
    \item Na UNIXu býval všechno soubor

    \pause

    \item Na soubor existuje souborový deskriptor

    \pause

    \item I na síťové spojení se vytváří souborový deskriptor
    \begin{itemize}
      \item Tzv.\ socket (\doc{man 7 socket})
    \end{itemize}

    \pause

    \item Na něm se dá provádět většina standardních operací
  \end{itemize}

\end{frame}

\begin{frame}[fragile]{\code{socket}}

  \begin{itemize}
    \item Otevření: \code{int socket(int domain, int type, int protocol)}
    \begin{itemize}
      \item Hlavičky: \header{sys/types.h}, \header{sys/socket.h}
      \item \code{domain}: protokol -- \code{AF_UNIX} \code{AF_INET},
	\code{AF_INET6}, \dots
      \item \code{type}: \code{SOCK_STREAM} (TCP), \code{SOCK_DGRAM} (UDP), \dots
      \item \code{protocol}: pokud je na výběr (NETLINK), jinak 0 (pro UDP i TCP)
      \item Návratová hodnota: souborový deskriptor nebo \code{\-1} při chybě
    \end{itemize}

    \pause

    \item Zavření
    \begin{itemize}
      \item Klasicky: \code{close}
\ifpcdirlinuxb
      \item Částečně: \code{int shutdown(int sockfd, int how)}
      \begin{itemize}
	\item \code{how} -- \code{SHUT_RD}, \code{SHUT_WR}, \code{SHUT_RDWR}
      \end{itemize}
\fi
    \end{itemize}

    \pause
  \end{itemize}

  \onslide<1->

  \begin{block}{Příklad (bez ošetření chyb)}
  \begin{lstlisting}
int s = socket(AF_UNIX, SOCK_STREAM, 0);
close(s);
  \end{lstlisting}
  \end{block}

\end{frame}

\fi

\begin{frame}{Úkol}

  \heading{Otevření TCP socketu}
  \begin{enumerate}
    \item Otevřete TCP socket na IPv4
    \begin{itemize}
      \item \code{socket}, \code{AF_INET} a \code{SOCK_STREAM}
    \end{itemize}
    \item Zavřete
    \item Přeložte, vyzkoušejte
    \item Zkopírujte si do dvou zdrojových souborů
    \begin{itemize}
      \item Nebo rozdělte kód v závislosti na příkazové řádce
      \begin{itemize}
	\item \code{if (argc > 1) server(); else client();}
      \end{itemize}
      \item Zárodek klienta a serveru
    \end{itemize}
  \end{enumerate}

\end{frame}

\begin{frame}[fragile]{Pořadí bytů a síť}

  \begin{block}{Číslo \code{0x12345678} reprezentované v paměti}
  \begin{center}
   \begin{tabular}{l|l|llll}
\multirow{2}{*}{\textbf{Pořadí}} & \textbf{Příklad} &
  \multicolumn{4}{c}{\textbf{Offset v paměti}} \\
              & \textbf{architektury}       &    0 &    1 &    2 &    3 \\\hline
Big-endian    & \texttt{arm}, \texttt{s390}, síť & 0x12 & 0x34 & 0x56 & 0x78 \\
Little-endian & \texttt{arm}, \xes          & 0x78 & 0x56 & 0x34 & 0x12 \\
Mixed-endian  & \texttt{PDP-11}             & 0x34 & 0x12 & 0x78 & 0x56 \\
  \end{tabular}
  \end{center}
  \end{block}

  \pause
  \bigskip

  \heading{Pro typy o velikosti větší než 1 je nutné znát pořadí}
  \begin{itemize}
    \item Interpretace dat z/do sítě
    \begin{itemize}
      \item Např.\ UDP i TCP porty jsou 16bitové
    \end{itemize}

    \pause

    \item Funkce (\header{arpa/inet.h})
    \begin{itemize}
      \item \code{ntohs}, \code{ntohl} -- ze sítě (\textit{Net TO Host})
      \item \code{htons}, \code{htonl} -- do sítě (\textit{Host TO Net})
      \item \code{*s} -- \code{short} (16 bitů), \code{*l} -- \code{long}
        (32 bitů)
    \end{itemize}
  \end{itemize}

\end{frame}

\begin{frame}{Adresování v Linuxu}

  \begin{block}{Příklady adres}
  \begin{center}
  \begin{tabular}{|l|l|l|l|} \hline
    \textbf{Protokol} & \textbf{Vel.} &
      \textbf{Struktura} & \textbf{Příklad adresy} \\ \hline
    Unix (\code{AF_UNIX})  & 108 & \code{sockaddr_un} & \code{/tmp/socket_xy} \\ \hline
    IPv4 (\code{AF_INET})  & 4   & \code{sockaddr_in} & \code{127.0.0.1} \\ \hline
    IPv6 (\code{AF_INET6}) & 16  & \code{sockaddr_in6} & \code{2a01:4240:2e27::1} \\ \hline
    IPX  (\code{AF_IPX})   & 4+6 & \code{sockaddr_ipx} & \code{fedcba98 1a2b3c5d7e9f} \\ \hline
  \end{tabular}
  \end{center}
  \end{block}

  \pause

  \begin{itemize}
    \item Každý protokol má svou strukturu pro adresu
    \begin{itemize}
      \item Dostatečně velká, aby uchovala adresu (a další)
    \end{itemize}

    \pause

    \item Funkce pracují s obecnou \code{struct sockaddr}
    \begin{itemize}
      \item Každá \code{sockaddr_*} má položku \code{family}
      \begin{itemize}
	\item Např.\ \code{sockaddr_in.sin_family = AF_INET} \\
	  \code{sockaddr_in6.sin6_family = AF_INET6}
      \end{itemize}

      \pause

      \item Funkce mají parametr velikost
      \begin{itemize}
	\item Např.\ \code{sizeof(struct sockaddr_in)}
      \end{itemize}
    \end{itemize}

    \pause

    \item Pojme jakoukoliv: \code{struct sockaddr_storage}
  \end{itemize}

\end{frame}

\begin{frame}[fragile]{Adresování v Linuxu -- shrnutí}

  \begin{center}
  \begin{tabular}{|l|l|} \hline
  \textbf{Typ adresy} & \multicolumn{1}{c|}{\textbf{Adresa v Linuxu}} \\ \hline
  Obecná & \begin{lstlisting}
struct sockaddr {
  sa_family_t sa_family; /* AF_INET, AF_INET6, ... */
  char sa_data[14];
};
struct sockaddr_storage; // large enough for any address
  \end{lstlisting}
  \\ \hline

  \uncover<2->{IPv4} &
  \begin{uncoverenv}<2->
  \begin{lstlisting}
struct sockaddr_in {
  sa_family_t    sin_family; /* address family: AF_INET */
  in_port_t      sin_port;   /* port in network byte order */
  struct in_addr sin_addr;   /* internet address */
};
struct in_addr {
  uint32_t       s_addr;     /* address in network byte order */
};
  \end{lstlisting}
  \end{uncoverenv}
  \\ \hline

  \uncover<3->{IPv6} &
  \begin{uncoverenv}<3->
  \begin{lstlisting}
struct sockaddr_in6 {
  sa_family_t     sin6_family;   /* AF_INET6 */
  in_port_t       sin6_port;     /* port number */
  struct in6_addr sin6_addr;     /* IPv6 address */
};
struct in6_addr {
  unsigned char   s6_addr[16];   /* IPv6 address */
};
  \end{lstlisting}
  \end{uncoverenv}
  \\ \hline
  \end{tabular}
  \end{center}

\end{frame}

\subsection*{Klient}
\frame{\subsectionpage}

\begin{frame}[fragile]{Získání adresy}

  \begin{itemize}
    \item \code{int getaddrinfo(const char *node, const char *service,} \\
      \code{const struct addrinfo *hints, struct addrinfo **res)}
    \begin{itemize}
      \item Hlavičky: \header{sys/types.h}, \header{sys/socket.h},
        \header{netdb.h}
      \item \code{node}: cílový stroj (adresa/jméno)
      \item \code{service}: cílový port (číslo/jméno)
      \item \code{hints}: omezení, co hledat (IPv4, IPv6 apod.), nebo NULL
      \item \code{res}: návratová hodnota
      \item Návratová hodnota: 0 je OK, jinak číslo chyby
      \begin{itemize}
	\item \code{const char *gai_strerror(int errcode)}
      \end{itemize}
    \end{itemize}

    \item<2-> Uvolnění: \code{void freeaddrinfo(struct addrinfo *res)}
  \end{itemize}

  \begin{uncoverenv}<1->
  \begin{block}{\code{struct addrinfo}}
  \begin{lstlisting}
int  ai_flags;              socklen_t        ai_addrlen;
int  ai_family;             struct sockaddr *ai_addr;
int  ai_socktype;           char            *ai_canonname;
int  ai_protocol;           struct addrinfo *ai_next;
  \end{lstlisting}
  \end{block}
  \end{uncoverenv}

\end{frame}

\begin{frame}{Úkol}

  \heading{Získání adres}
  \begin{enumerate}
    \item Doplňte do klienta z předchozího příkladu za vytvoření socketu
    \item Použijte \code{getaddrinfo}

    \item Zjistěte, kolik je k dispozici adres pro server \cmd{google.com}
    \begin{itemize}
      \item Parametr \code{node}
    \end{itemize}

    \item Omezte se jen na TCP a port HTTP
    \begin{itemize}
      \item \code{hints.ai_socktype} bude \code{SOCK_STREAM}
      \item \code{hints.ai_protocol} bude \code{IPPROTO_TCP}
      \item \code{service} bude \code{\"http\"}
    \end{itemize}

    \item Přeložte, vyzkoušejte
  \end{enumerate}

\end{frame}

\begin{frame}[fragile]{Převod adresy na řetězec}

  \begin{block}{Prototyp}
  \begin{lstlisting}[aboveskip=2pt,belowskip=3pt]
int getnameinfo(const struct sockaddr *addr, socklen_t addrlen, char *host,
  socklen_t hostlen, char *serv, socklen_t servlen, int flags)
  \end{lstlisting}
  \end{block}

  \begin{itemize}
    \item Hlavičky: \header{sys/socket.h}, \header{netdb.h}
    \item \code{addr}: adresa (např.\ \code{ai.ai_addr})
    \item \code{addrlen}: velikost adresy (např.\ \code{ai.ai_addrlen})
    \item \code{host}: buffer pro jméno/adresu (vel.\ \code{NI_MAXHOST})
    \item \code{serv}: buffer pro službu (vel.\ \code{NI_MAXSERV})
    \item \code{hostlen}/\code{servlen}: velikost bufferu \code{host} a
      \code{serv}
    \item \code{flags}: \code{NI_NUMERICHOST}, \code{NI_NUMERICSERV}, \dots

    \item Návratová hodnota: 0 je OK, jinak číslo chyby
    \begin{itemize}
      \item \code{const char *gai_strerror(int errcode)}
    \end{itemize}
  \end{itemize}

\end{frame}

\begin{frame}{Úkol}

  \heading{Výpis adres}
  \begin{enumerate}
    \item Doplňujte do předchozího příkladu
    \item Použijte \code{getnameinfo}

    \item Vypište i adresy
    \begin{itemize}
      \item \code{serv} může být NULL
      \item Nastavte \code{flags} na \code{NI_NUMERICHOST}
    \end{itemize}

    \item Přeložte, vyzkoušejte
  \end{enumerate}

\end{frame}

\begin{frame}[fragile]{Klient}

  \begin{itemize}
    \item Připojení: { \footnotesize
      \code{int connect(int sockfd, const struct sockaddr *addr,} \\
      \code{socklen_t addrlen)}
      }
    \begin{itemize}
      \item \code{sockfd}: návratová hodnota funkce \code{socket}
      \item \code{addr}: adresa, kam se připojit
      \item \code{addrlen}: velikost struktury adresy
    \end{itemize}

\ifpcdirlinux
    \pause

    \item Příjem: \code{read} nebo \code{ssize_t recv(int sockfd, void *buf,
      size_t len, int flags)}

    \item Vysílání: \code{write} nebo \code{ssize_t send(int sockfd, const void
      *buf, size_t len, int flags)}
\fi
  \end{itemize}

  \pause

  \begin{block}{Příklad (bez ošetření chyb)}
  \begin{lstlisting}
getaddrinfo("nekde", "neco", &hint, &ai);
fd = connect(sockfd, ai->ai_addr, ai->ai_addrlen);
write(fd, "ahoj", 4);
  \end{lstlisting}
  \end{block}

\end{frame}

\begin{frame}{Úkol}

  \heading{Klient}
  \begin{enumerate}
    \item Připojte se na \code{localhost} na port \code{ssh}
    \item Přečtěte ze socketu 20 znaků
    \item Vypište je na standardní výstup
    \item Vyzkoušejte
    \item Co vám ssh server poslal?
  \end{enumerate}

\end{frame}

\subsection*{Server}
\frame{\subsectionpage}

\begin{frame}{Server}

  \heading{Tři kroky}
  \begin{enumerate}
    \item Nastavení adresy a portu
    \begin{itemize}
      \item \code{int bind(int sockfd, const struct sockaddr *addr,}\\
        \code{socklen_t addrlen)}
      \begin{itemize}
	\item \code{addr}: IP adresa se nastavuje často na \code{INADDR_ANY}
	  (IPv4, host) nebo \code{in6addr_any}/\code{IN6ADDR_ANY_INIT} (IPv6,
	  net)
      \end{itemize}
    \end{itemize}

    \pause

    \item Nastavení socketu jako pasivního (bude čekat na spojení)
    \begin{itemize}
      \item \code{int listen(int sockfd, int backlog)}
      \begin{itemize}
	\item \code{backlog}: kolik spojení může čekat na obsluhu
      \end{itemize}

      \item Od tohoto okamžiku bude fungovat \code{connect} na tento port
    \end{itemize}

    \pause

    \item Čekání na příchozí spojení
    \begin{itemize}
      \item \code{int accept(int sockfd, struct sockaddr *addr, } \\
        \code{socklen_t *addrlen)}
      \begin{itemize}
	\item \code{addr}: návratová hodnota -- adresa toho, kdo se připojil
      \end{itemize}

      \item Volá se opakovaně pro obsluhu dalších spojení
    \end{itemize}

  \end{enumerate}

\end{frame}

\begin{frame}{Úkol}

  \heading{Server -- echo}
  \begin{enumerate}
    \item Vytvořte server
    \item Pro každé příchozí spojení vytvořte potomka
    \item Rodič bude čekat na další spojení
    \item Potomek obslouží spojení
    \item Přečte ze socketu až 100 znaků
    \item Pošle je zpět
    \item Upravte klienta, aby po připojení něco poslal
    \item Potom přečte ze socketu až 100 znaků
    \item Vypíše je
    \item Vyzkoušejte
  \end{enumerate}

\end{frame}

\begin{frame}{Multiplexing}

  \heading{Jak přepínat mezi čtením, zápisem a různými deskriptory}
  \begin{itemize}
    \item \code{int poll(struct pollfd *fds, nfds_t nfds, int timeout)} (\header{poll.h})
    \begin{itemize}
      \item \code{fds}: pole popisující, na které \code{fd} čekat (a jaké
        události)
      \item \code{nfds}: velikost pole \code{fds}
      \item \code{timeout}: \code{NULL} nebo timeout
      \item \code{fds} je po návratu patřičně upraveno
    \end{itemize}

    \pause

    \item \code{int select(int nfds, fd_set *readfds, fd_set *writefds,
      fd_set *exceptfds, struct timeval *timeout)} (\header{sys/select.h})
    \begin{itemize}
      \item \code{nfds}: nejvyšší číslo \code{fd} v množinách + 1
      \item \code{readfds}: čekat, až je bude možno číst, nebo \code{NULL}
      \item \code{writefds}: čekat, až do nich bude možno zapisovat, nebo
	\code{NULL}
      \item \code{exceptfds}: čekat, až nastane chyba, nebo \code{NULL}
      \item Návratová hodnota: počet nastavených bitů v \code{readfds},
	\code{writefds} a \code{exceptfds}
    \end{itemize}

  \end{itemize}

\end{frame}

\begin{frame}{Úkol}

  \heading{Multiplexing}
  \begin{enumerate}
    \item Upravte klienta
    \item Bude číst zároveň ze spojení a standardního vstupu
    \item Data ze spojení opět vypište
    \item Data ze standardního vstupu pošlete serveru
    \item Vyzkoušejte
  \end{enumerate}

\end{frame}

